\documentclass[12pt]{article}

\usepackage[spanish]{babel}
\usepackage[utf8]{inputenc}
\usepackage[T1]{fontenc}
\usepackage{amsmath, amssymb}
\usepackage{geometry}
\usepackage{xcolor}

\geometry{letterpaper, margin=2cm}

\begin{document}
	
	\section*{Solución del problema}
	
	\textbf{Datos:}
	
	\begin{itemize}
		\item Dosis recomendada: $140 \text{ kg/ha}$
		\item Superficie: $3{,}6 \text{ ha}$
	\end{itemize}
	
	\vspace{0.3cm}
	
	\textbf{1. Verificación del cálculo}
	
	El técnico realiza:
	\[
	140 \times 3{,}6 = 504
	\]
	
	Descomponiendo:
	\[
	140 \times 3 = 420 \quad \text{y} \quad 140 \times 0{,}6 = 84
	\]
	
	\[
	420 + 84 = 504
	\]
	
	\textbf{Conclusión:} El cálculo es \textbf{correcto}. El resultado es $504 \text{ kg}$.
	
	\vspace{0.3cm}
	
	\textbf{2. Análisis de razonabilidad}
	
	Si:
	\[
	1 \text{ ha} \rightarrow 140 \text{ kg}
	\]
	
	Entonces:
	\[
	3{,}6 \text{ ha} \rightarrow 140 \times 3{,}6 = 504 \text{ kg}
	\]
	
	Además, como $3{,}6 \approx 4$:
	\[
	140 \times 4 = 560 \text{ kg}
	\]
	
	El valor obtenido ($504$ kg) es menor que $560$ kg, lo cual es coherente ya que $3{,}6 < 4$.
	
	\textbf{Conclusión:} El resultado es \textbf{razonable}.
	
	\vspace{0.3cm}
	
	\textbf{3. Análisis del resultado}
	
	\begin{itemize}
		\item \textbf{Unidad:}
		
		\[
		\text{kg/ha} \times \text{ha} = \text{kg}
		\]
		
		La unidad final es \textbf{kilogramos (kg)}.
		
		\item \textbf{Magnitud:}
		
		\[
		\frac{504}{3{,}6} = 140 \text{ kg/ha}
		\]
		
		El resultado coincide con la dosis recomendada.
		
		\item \textbf{Coherencia del contexto:}
		
		El cálculo es proporcional y consistente con la recomendación técnica.
	\end{itemize}
	
	\vspace{0.3cm}
	
	\textbf{4. Cálculo de sacos}
	
	Cada saco contiene $25$ kg:
	
	\[
	\frac{504}{25} = 20{,}16
	\]
	
	\textbf{Cantidad de sacos:}
	
	No es posible comprar fracciones de saco, por lo tanto:
	
	\[
	\boxed{21 \text{ sacos}}
	\]
	
	\textbf{Justificación del redondeo:}
	
	\begin{itemize}
		\item $20$ sacos $\rightarrow 20 \times 25 = 500$ kg (insuficiente)
		\item $21$ sacos $\rightarrow 21 \times 25 = 525$ kg (suficiente)
	\end{itemize}
	
	Se debe redondear \textbf{hacia arriba} para asegurar la cantidad necesaria de fertilizante.
	
	\vspace{0.3cm}
	
	\textbf{Conclusión final:}
	
	\begin{itemize}
		\item El cálculo del técnico es correcto.
		\item El resultado es razonable y coherente.
		\item La unidad es kg.
		\item Se deben comprar \textbf{21 sacos}.
	\end{itemize}
	
\end{document}